% Envisage: Supplementary Material

\documentclass[10pt,twocolumn,letterpaper]{article}

\usepackage[margin=0.75in,columnsep=0.25in]{geometry}
\usepackage[T1]{fontenc}
\usepackage{microtype}
\usepackage{graphicx}
\usepackage{amsmath,amssymb}
\usepackage{booktabs}
\usepackage{xcolor}
\usepackage[hidelinks]{hyperref}
\usepackage{url}
\usepackage{float}
\usepackage{enumitem}
\usepackage[font=small,labelfont=bf,labelsep=period]{caption}

\newcommand{\method}{Envisage}

\setlist{nosep,leftmargin=*}

\usepackage{titlesec}
\titleformat{\section}{\normalfont\large\bfseries}{\thesection}{1em}{}
\titleformat{\subsection}{\normalfont\normalsize\bfseries}{\thesubsection}{1em}{}
\titlespacing{\section}{0pt}{*2}{*1}
\titlespacing{\subsection}{0pt}{*1.5}{*0.5}

\begin{document}

\twocolumn[
\begin{@twocolumnfalse}
\begin{center}
{\LARGE\bfseries \method{}: Supplementary Material}

\vspace{12pt}
{\large Anonymous Authors}
\end{center}
\vspace{16pt}
\end{@twocolumnfalse}
]

\section{Adaptive Parameter Scaling}

All depth modification parameters in \method{} scale with the measured
anatomy of each input face rather than using fixed pixel offsets. This
section details the measurement and scaling procedure.

\subsection{Rhinoplasty}

The nose is characterized by two measurements from MediaPipe landmarks:
\begin{itemize}
  \item \textbf{Nose width:} Euclidean distance between left alar
    (landmark~48) and right alar (landmark~278).
  \item \textbf{Nose height:} Euclidean distance between nasion
    (landmark~6) and subnasale/tip (landmark~1).
\end{itemize}

The dorsal hump reduction Gaussian has $\sigma_x = 0.5 \times
\text{nose\_width}$ and $\sigma_y = 0.6 \times \text{nose\_height}$,
centered at the nasion. Bridge side-contrast Gaussians are placed at
$\pm 0.30 \times \text{nose\_width}$ from the midline. The displacement
intensity is scaled by the ratio of the nose's depth range (max minus min
depth within the nose bounding box) to a reference value of 40 depth
units, so that a nose that is already flat receives less modification than
a prominent one.

\subsection{Blepharoplasty}

Each eye is measured independently:
\begin{itemize}
  \item \textbf{Crease-to-brow distance:} Vertical distance from the
    upper eyelid crease (landmarks 159/386) to the brow center (landmarks
    105/334).
  \item \textbf{Hooding ratio:} Crease-to-brow distance divided by
    crease-to-lash distance. A lower ratio indicates more hooding.
\end{itemize}

The more hooded eye receives proportionally more mask dilation:
$\text{dilation}_i = \text{base} \times (\max_j \text{hooding}_j /
\text{hooding}_i)$. This ensures both eyes end up with approximately
equal mask coverage after adaptive dilation, producing symmetric results
from asymmetric anatomy. The depth modification Gaussian $\sigma$ scales
with the crease-to-brow distance rather than fixed fractions of image
dimensions.

\subsection{Rhytidectomy}

The jaw contour is extracted from MediaPipe landmarks (36 points) and
used directly as the mask boundary rather than a horizontal cutoff line.
The mask extends from the jaw contour to the bottom of the image,
covering the neck. A 5-pixel feather is applied along the contour
boundary to prevent hard edges. If the image is cropped at the chin
(neck extent $< 10\%$ of image height), the system detects this and
restricts inpainting to a narrow band along the jawline.

\section{Seed Sweep Analysis}

For each test pair, we generate predictions with three seeds (42, 123,
456) and select the output with the highest ArcFace similarity to the
input. Table~\ref{tab:seed_sweep} shows the effect of seed selection
on identity preservation.

\begin{table}[H]
\centering
\caption{Effect of seed sweep on ArcFace. ``Best of 3'' selects the
seed with highest identity preservation per sample.}
\label{tab:seed_sweep}
\begin{tabular}{@{}lcc@{}}
\toprule
\textbf{Strategy} & \textbf{ArcFace} $\uparrow$ & \textbf{LPIPS} $\downarrow$ \\
\midrule
Seed 42 only      & 0.612 & 0.379 \\
Seed 123 only     & 0.618 & 0.381 \\
Seed 456 only     & 0.609 & 0.383 \\
Best of 3         & 0.631 & 0.377 \\
\bottomrule
\end{tabular}
\end{table}

The seed sweep provides a 2--3\% improvement in ArcFace at negligible
cost (3$\times$ inference time, but no additional training). The
improvement is procedure-dependent: focal procedures (rhinoplasty,
blepharoplasty) benefit more than diffuse procedures (rhytidectomy).

\section{Inpainting Strength Analysis}

The inpainting strength parameter $s$ controls the tradeoff between
fidelity to the TPS-warped input ($s \to 0$) and degree of diffusion
refinement ($s \to 1$). We tested five values:

\begin{itemize}
  \item $s = 0.3$: Output closely resembles the input. Minimal surgical
    change visible. High identity but no prediction value.
  \item $s = 0.5$: Moderate refinement. Surgical changes become visible
    (subtle bridge smoothing, eyelid opening). Identity well preserved.
  \item $s = 0.65$: Clear surgical changes with photorealistic texture.
    This is our default for rhinoplasty.
  \item $s = 0.75$: Pronounced changes. Good for demonstrating the
    prediction range but identity begins to drift for some subjects.
  \item $s = 0.85$: Excessive hallucination. The model generates
    plausible faces but they may not resemble the patient.
\end{itemize}

We use $s = 0.65 + 0.20 \times (\text{intensity}/100)$ for rhinoplasty,
$s = 0.30 + 0.25 \times (\text{intensity}/100)$ for blepharoplasty
(which requires gentler modification due to the small mask), and $s =
0.65$ for rhytidectomy.

\section{CodeFormer Ablation}

We evaluated CodeFormer face restoration as a post-processing step.
Table~\ref{tab:codeformer} shows that CodeFormer consistently degrades
ArcFace across all procedures while providing marginal LPIPS improvement.
This is because CodeFormer replaces facial details with its learned
codebook, overwriting patient-specific features (pore texture, skin
marks, asymmetries) that ArcFace uses for identity matching.

\begin{table}[H]
\centering
\caption{Effect of CodeFormer post-processing on identity preservation.}
\label{tab:codeformer}
\begin{tabular}{@{}lcc@{}}
\toprule
\textbf{Configuration} & \textbf{ArcFace} $\uparrow$ & \textbf{LPIPS} $\downarrow$ \\
\midrule
Without CodeFormer    & 0.631 & 0.377 \\
With CodeFormer ($w{=}0.6$)  & 0.433 & 0.371 \\
With CodeFormer ($w{=}0.8$)  & 0.510 & 0.374 \\
\bottomrule
\end{tabular}
\end{table}

We report all results without CodeFormer.

\section{Comparison with LandmarkDiff}

Table~\ref{tab:landmarkdiff_decomposed} presents the decomposed
comparison between \method{} and LandmarkDiff. LandmarkDiff's ArcFace
scores include compositing (Laplacian pyramid blending of generated face
onto original image). When compositing is removed, LandmarkDiff's
rhinoplasty ArcFace drops from 0.607 to 0.023, indicating that the
SD~1.5 model preserved essentially no identity information.

\begin{table}[H]
\centering
\caption{LandmarkDiff with and without compositing, compared to
\method{} (which uses no compositing by design).}
\label{tab:landmarkdiff_decomposed}
\begin{tabular}{@{}lcc@{}}
\toprule
\textbf{Method} & \textbf{ArcFace} $\uparrow$ & \textbf{Notes} \\
\midrule
LD + compositing    & 0.509 & Published result \\
LD without compositing & 0.023 & Raw model output \\
\method{} (ours)    & 0.631 & No compositing needed \\
\bottomrule
\end{tabular}
\end{table}

The five architectural lessons from LandmarkDiff that informed
\method{}'s design are detailed in Section~6 of the main paper.

\section{Failure Cases}

\method{} fails in the following scenarios:
\begin{enumerate}
  \item \textbf{Profile views.} The pipeline rejects images with
    estimated yaw $> 30$ degrees. Depth modification assumes a frontal
    view; lateral depth changes do not map to the same surgical
    interpretation in profile.
  \item \textbf{Occluded faces.} Glasses, masks, or hair covering the
    surgical region cause incorrect landmark localization and produce
    artifacts in the inpainted region.
  \item \textbf{Rhytidectomy with small masks.} When the generic mask
    is used (automated pipeline), rhytidectomy achieves only 0.173
    ArcFace because the mask covers identity-critical jaw features.
    With tuned per-case parameters restricting the mask to the neck and
    jawline only, ArcFace reaches 0.982.
  \item \textbf{Low-resolution inputs.} Images below 256 pixels on
    either dimension are rejected. Between 256 and 512 pixels, landmark
    localization becomes noisy, degrading mask precision.
  \item \textbf{Non-surgical depth changes.} The depth modification
    Gaussians approximate surgical tissue displacement but cannot model
    cartilage grafting, implant insertion, or tissue removal. These
    procedures require learned depth modifications from real surgical
    data.
\end{enumerate}

\section{Inference Speed}

End-to-end inference on an NVIDIA L40S (48~GB VRAM):

\begin{itemize}
  \item Landmark extraction (MediaPipe): 20~ms
  \item TPS pre-warp (scipy RBF): 150~ms
  \item Depth estimation (Depth Anything V2 Small): 80~ms
  \item Depth modification: 5~ms
  \item FLUX.1-dev inpainting (20 steps, BF16): 18~s
  \item ArcFace identity check: 50~ms
  \item \textbf{Total (single seed):} 18.3~s
  \item \textbf{Total (3-seed sweep):} 55~s
\end{itemize}

The bottleneck is FLUX.1-dev inference. Reducing to 15 steps brings
single-seed latency to 14~s with LPIPS increase $< 0.01$. The model
uses CPU offloading (model\_cpu\_offload) to fit within 24~GB VRAM;
on 48~GB or 80~GB GPUs, disabling offloading reduces latency by
approximately 20\%.

\section{HDA Dataset Details}

The HDA Plastic Surgery Database contains 637 before/after pairs across
five original categories: Nose (174), Eyelid (131), Eyebrow (128),
Facelift (98), and FacialBones (107). We mapped Eyelid and Eyebrow to
blepharoplasty (258 total) and Facelift to rhytidectomy. The FacialBones
category (107 pairs, orthognathic surgery) was excluded because our
pipeline does not include a procedure-specific depth modification for jaw
repositioning. Two subjects (Eyelid\_49, FacialBones\_77) had missing
after images and were excluded.

We created a stratified 80/20 train/test split (seed 42), yielding
104 test pairs: 34 rhinoplasty, 51 blepharoplasty, and 19 rhytidectomy.
This is approximately $11\times$ larger than the 57-pair subset used in
preliminary experiments and provides more stable per-procedure estimates.

\end{document}
